\section{Building Blocks}
\label{sec:blocks}

Let us revisit Eqs.~\eqref{eq:plonk1} and \eqref{eq:plonk2}.
Our goal is to prove both relations using bitwidth-aware protocols.
Eq.~\eqref{eq:plonk1} is a zero check. By Corollary~\ref{corollary:schwartz}, it suffices to project its
left-hand side onto a random tensor-structured vector
\(
    \ext{\bm{\xi}}
    =
    \tensorspan{\ext{\bm z}}
\)
and prove that the resulting inner product is zero. This reduces
the zero check to proving an inner-product claim for each term in
Eq.~\eqref{eq:plonk1}. Eq.~\eqref{eq:plonk2}, in turn, can
be established using a permutation argument. In this section, we
develop bitwidth-aware versions of these protocol components.

\subsection{Packed Inner-Product Proof}
\label{sec:ip}

We first present a bitwidth-aware inner-product proof that packs
\(\kappa\) base-field blocks into a single extension-field vector.
It reduces a length-\(N\) base-field inner product to a
length-\(m\) extension-field inner product, where \(N=m\kappa\).
The same technique will later be used to construct bitwidth-aware
Hadamard-product and permutation proofs. This protocol plays the
same role of the degree-one zero-check in~\cite{iacr/DaoT24}.

Consider the relation
\[
    c=\langle\bm a,\bm b\rangle,
    ~
    \bm a,\bm b\in\Field^N.
\]
Partition each vector into \(\kappa\) blocks of length \(m\), and
arrange these blocks as columns:
\begin{equation}
\label{eq:matAB}
    \bm A
    \coloneqq
    \concat{
        \bm a^{(1)},\ldots,\bm a^{(\kappa)}
    },
    \qquad
    \bm B
    \coloneqq
    \concat{
        \bm b^{(1)},\ldots,\bm b^{(\kappa)}
    }
    \in\Field^{m\times\kappa}.
\end{equation}
Then, we encode the blocks as the vector-valued polynomials
\[
    \bm g_A(X)
    \coloneqq
    \sum_{j=1}^{\kappa}\bm a^{(j)}X^{j-1},
    \qquad
    \bm g_B(X)
    \coloneqq
    \sum_{j=1}^{\kappa}\bm b^{(j)}X^{j-1}
    \in\Field[X]^m,
\]
and define the Laurent polynomial
\begin{equation}
    f(X)
    \coloneqq
    \left\langle
        \bm g_A(X),\bm g_B(X^{-1})
    \right\rangle.
\label{eq:ip-laurent-polynomial}
\end{equation}
The key observation is that
\(c=\langle\bm a,\bm b\rangle\) if and only if there exist
\(h_+(X),h_-(X)\in\Field[X]\), each of degree at most
\(\kappa-2\), such that
\begin{equation}
    f(X)
    =
    c
    +Xh_+(X)
    +X^{-1}h_-(X^{-1}).
\label{eq:ip-correction-identity}
\end{equation}

After the commitments to \(\bm A\) and \(\bm B\) are fixed, the
prover sends \(h_+\) and \(h_-\). The verifier then samples
\(
    \ext{x}\sample\ExtensionField^\times.
\)
Define the packed vectors
\[
    \ext{\bm a}
    \coloneqq
    \bm g_A(\ext{x}),
    \qquad
    \ext{\bm b}
    \coloneqq
    \bm g_B(\ext{x}^{-1})
    \in\ExtensionField^m.
\]
Evaluating Eq.~\eqref{eq:ip-correction-identity} at \(\ext{x}\)
reduces the original claim to
\begin{equation}
    \left\langle
        \ext{\bm a},\ext{\bm b}
    \right\rangle
    =
    c
    +\ext{x}h_+(\ext{x})
    +\ext{x}^{-1}h_-(\ext{x}^{-1}).
\label{eq:packed-inner-product-check}
\end{equation}
The verifier computes the right-hand side directly, while the
left-hand side is proved using a length-\(m\) inner-product proof
over \(\ExtensionField\). The complete protocol is given in
Algorithm~\ref{alg:ip} in Appendix~\ref{appx:pseudocode}.

\begin{proposition}[Packed inner-product reduction]
\label{proposition:ip-complexity}
Algorithm~\ref{alg:ip} is complete and has soundness error at most
\[
    \frac{2\kappa-2 }{|\ExtensionField|-1} + \frac{2\log_2 m}{|\ExtensionField|}.
\]
Excluding the terminal MLE openings, its prover reduction cost is
\[
    N(\kappa-1)\,\bb
    +
    2N\,\be
    +
    \left(
        \frac{4N}{\kappa}
        +
        O\left(
            \kappa+\log\frac{N}{\kappa}
        \right)
    \right)\ee.
\]
Its peak prover memory is \(O(N+\kappa^2)\) base-field
coordinates, which is \(O(N)\) whenever \(m\geq\kappa\).
\end{proposition}

The proof is deferred to Appendix~\ref{appx:proof-ip}.

\subsection{Packed Hadamard-Product Proof}
\label{sec:hadamard}

We next extend the packing technique to the weighted
Hadamard-product relation
\begin{equation}
    \ext{c}
    =
    \left\langle
        \ext{\bm\xi},
        \bm a\circ\bm b
    \right\rangle,
    \qquad
    \ext{\bm\xi}
    =
    \ext{\bm r}\kronecker\ext{\bm l},
\label{eq:hadamard-claim}
\end{equation}
where
\[
    \bm a,\bm b\in\Field^N,
    \qquad
    \ext{\bm l}\in\ExtensionField^m,
    \qquad
    \ext{\bm r}\in\ExtensionField^\kappa,
    \qquad
    N=m\kappa.
\]
Using the block matrices \(\bm A,\bm B\in\Field^{m\times\kappa}\)
defined in Eq.~\eqref{eq:matAB}, Eq.~\eqref{eq:hadamard-claim}
can be written as
\begin{equation*}
    \ext{c}
    =
    \ext{\bm l}^{\top}
    (\bm A\circ\bm B)
    \ext{\bm r}.
\end{equation*}

The prover first computes the column claims
\[
    \ext{\bm c}^{\top}
    \coloneqq
    \ext{\bm l}^{\top}
    (\bm A\circ\bm B)
    \in\ExtensionField^\kappa,
\]
or equivalently,
\[
    \ext{c}_j
    =
    \left\langle
        \ext{\bm l},
        \bm a^{(j)}\circ\bm b^{(j)}
    \right\rangle
    \qquad
    \text{for every }j\in[\kappa].
\]
The verifier checks
\begin{equation*}
    \ext{c}
    \EqOrNot
    \left\langle
        \ext{\bm c},\ext{\bm r}
    \right\rangle.
\end{equation*}
It remains to prove that all entries of \(\ext{\bm c}\) are
correct.

Let \(\{\psi_j(X)\}_{j=1}^{\kappa}\) and \(\Psi(X)\) be the
Lagrange basis and vanishing polynomial from
Definition~\ref{def:lagrange}. Define
\[
\begin{aligned}
    \bm g_A(X)
    &\coloneqq
    \sum_{j=1}^{\kappa}
    \bm a^{(j)}\psi_j(X),
    \quad
    \bm g_B(X)
    \coloneqq
    \sum_{j=1}^{\kappa}
    \bm b^{(j)}\psi_j(X),
    \quad
    g_C(X)
    \coloneqq
    \sum_{j=1}^{\kappa}
    \ext{c}_j\psi_j(X).
\end{aligned}
\]
Here,
\[
    \bm g_A(X),\bm g_B(X)\in\Field[X]^m,
    \qquad
    g_C(X)\in\ExtensionField[X].
\]
By the defining property of the Lagrange basis,
\[
    \bm g_A(\alpha_j)=\bm a^{(j)},
    \qquad
    \bm g_B(\alpha_j)=\bm b^{(j)},
    \qquad
    g_C(\alpha_j)=\ext{c}_j.
\]
Consequently, all column claims are correct if and only if
\[
    f(\alpha_j)=0
    \qquad
    \text{for every }j\in[\kappa],
\]
where
\begin{equation*}
    f(X)
    \coloneqq
    \left\langle
        \ext{\bm l},
        \bm g_A(X)\circ\bm g_B(X)
    \right\rangle
    -
    g_C(X).
\end{equation*}
Since \(\deg f\leq 2\kappa-2\), this is equivalent to the
existence of a polynomial
\[
    h(X)\in\ExtensionField[X],
    \qquad
    \deg h\leq\kappa-2,
\]
such that
\begin{equation}
    f(X)=\Psi(X)h(X).
\label{eq:hadamard-quotient}
\end{equation}

The prover sends \(\ext{\bm c}\) and \(h(X)\) before the verifier
samples
\[
    \ext{x}
    \sample
    \ExtensionField\setminus\mathcal H.
\]
The verifier evaluates the Lagrange basis at \(\ext{x}\) and sets
\[
    \ext{\bm\rho}(\ext{x})
    \coloneqq
    \bigl(
        \psi_1(\ext{x}),\ldots,\psi_\kappa(\ext{x})
    \bigr)^\top.
\]
The corresponding packed vectors are
\[
\begin{aligned}
    \ext{\bm a}
    &\coloneqq
    \bm A\ext{\bm\rho}(\ext{x})
    =
    \bm g_A(\ext{x}),
    \qquad
    \ext{\bm b}
    \coloneqq
    \bm B\ext{\bm\rho}(\ext{x})
    =
    \bm g_B(\ext{x}),
\end{aligned}
\]
and the verifier computes
\[
    \ext{c}_x
    \coloneqq
    g_C(\ext{x})
    =
    \left\langle
        \ext{\bm c},
        \ext{\bm\rho}(\ext{x})
    \right\rangle.
\]
Evaluating Eq.~\eqref{eq:hadamard-quotient} at \(\ext{x}\)
reduces the Hadamard-product relation to
\begin{equation*}
    \left\langle
        \ext{\bm a},
        \ext{\bm b}\circ\ext{\bm l}
    \right\rangle
    =
    \ext{c}_x
    +
    \Psi(\ext{x})h(\ext{x}).
\end{equation*}
This is a single length-\(m\) inner-product relation over
\(\ExtensionField\), which is reduced using
Fact~\ref{fact:ip}.

Let
\[
    \ext{\bm y}\in\ExtensionField^{\log_2m},
    \qquad
    \ext{v}_1,\ext{v}_2\in\ExtensionField
\]
be the challenge point and terminal scalars returned by that
reduction. Its terminal claims can be expressed directly against
the original vectors as
\begin{align}
\label{eq:hadamard-terminal}
    \ext{v}_1
    &\EqOrNot
    \left\langle
        \bm a,
        \ext{\bm\rho}(\ext{x})
        \kronecker
        \tensorspan{\ext{\bm y}}
    \right\rangle,
    \qquad
    \ext{v}_2
    \EqOrNot
    \left\langle
        \bm b,
        \ext{\bm\rho}(\ext{x})
        \kronecker
        \left(
            \ext{\bm l}
            \circ
            \tensorspan{\ext{\bm y}}
        \right)
    \right\rangle.
\end{align}

Algorithm~\ref{alg:hadamard} in Appendix~\ref{appx:pseudocode}
summarizes the protocol.

\begin{proposition}[Packed Hadamard-product proof]
\label{proposition:hadamard-complexity}
Assume that the terminal MLE claims output by
Algorithm~\ref{alg:hadamard} are subsequently verified. Then the
packed Hadamard-product proof is complete and has soundness error
at most
\[
\begin{aligned}
    \epsilon_{\mathsf{HP}}
    &\leq
    \frac{2\kappa-2}
         {|\ExtensionField|-\kappa}
    +
    \frac{2\log_2N}
         {|\ExtensionField|}.
\end{aligned}
\]
Assuming that the coefficients of the Lagrange basis and
\(\Psi(X)\) are precomputed, its dominant prover cost under
schoolbook arithmetic is
\[
    3\kappa N\,\bb
    +
    7N\,\be
    +
    \frac{6N}{\kappa}\,\ee.
\]
More precisely, the prover cost is
\[
\begin{aligned}
    &(3\kappa+1)N\,\bb+
    \left(
        7N-\frac{N}{\kappa}
        +O(\kappa^2)
    \right)\be+
    \left(
        \frac{6N}{\kappa}
        +
        O\left(
            \kappa+\log\frac{N}{\kappa}
        \right)
    \right)\ee.
\end{aligned}
\]
The peak prover memory is \(O(N+\kappa^2)\) base-field
coordinates, which is \(O(N)\) whenever \(m\geq\kappa\).
\end{proposition}

The proof is deferred to Appendix~\ref{appx:proof-hadamard}.

\subsection{Packed Permutation Proof}
\label{sec:permutation}

We then apply the packed Hadamard-product proof to the permutation
relation:
\[
    \bm a\stackrel{\mathsf{perm}}{=}\bm b,
    \qquad
    \bm a,\bm b\in\Field^N.
\]
After the commitments to \(\bm a\) and \(\bm b\) are fixed, the
verifier samples
\(
    \ext{z}\sample\ExtensionField
\)
and reduces the permutation relation to the fingerprint check
\begin{equation}
   \Theta \coloneqq  \mathsf{GrandProd}
    \bigl(\ext{z}\bm 1+\bm a\bigr)
    \EqOrNot
    \mathsf{GrandProd}
    \bigl(\ext{z}\bm 1+\bm b\bigr).
\label{eq:packed-permutation-fingerprint}
\end{equation}
If \(\bm a\) is a permutation of \(\bm b\), this identity holds
for every \(\ext z\). Conversely, if the two vectors are not
permutations, then
\[
    \prod_{i=1}^{N}(Z+a_i)
    -
    \prod_{i=1}^{N}(Z+b_i)
\]
is a nonzero polynomial of degree at most \(N-1\). Therefore,
Eq.~\eqref{eq:packed-permutation-fingerprint} holds for a
uniformly random \(\ext z\) with probability at most
\(
    \min\left\{
        1,\frac{N-1}{|\ExtensionField|}
    \right\}.
\)

We next give a bitwidth-aware proof of one side of
Eq.~\eqref{eq:packed-permutation-fingerprint}. The same proof is
applied to both vectors when \(\bm a\) and \(\bm b\) are
committed. If \(\bm b\) is public, its grand product can instead
be computed directly.

Let \(\bm d\in\Field^N\) denote the vector whose grand product is
to be proved. Assume that
\[
    N=m\kappa,
    \qquad
    m,\kappa\in2^{\mathbb N},
\]
and arrange its \(\kappa\) length-\(m\) blocks as
\[
    \bm D
    \coloneqq
    \concat{
        \bm d^{(1)},\ldots,\bm d^{(\kappa)}
    }
    \in\Field^{m\times\kappa}.
\]
Identical padding may be applied when these divisibility
conditions do not hold.

Let \(\log_2\kappa\). Define the formal column-product
tree
\begin{equation}
\begin{aligned}
    \bm C_0(Z)
    &\coloneqq
    Z\bm 1_{m\times\kappa}+\bm D,\\
    \bm C_t(Z)
    &=
    \concat{
        \bm C_{t,L}(Z),
        \bm C_{t,R}(Z)
    },\\
    \bm C_{t+1}(Z)
    &\coloneqq
    \bm C_{t,L}(Z)
    \circ
    \bm C_{t,R}(Z),
    \qquad
    t\in\{0,\ldots,\log_2\kappa-1\}.
\end{aligned}
\label{eq:packed-permutation-product-tree}
\end{equation}
At level \(t\),
\[
    \bm C_t(Z)
    \in
    \Field[Z]^{m\times(\kappa/2^t)}
    \qquad\text{and}\qquad
    \deg_Z\bm C_t\leq2^t.
\]
The final layer is a single column
\[
    \bm C_{\log_2 \kappa}(Z)
    \in\Field[Z]^m
\]
satisfying
\begin{equation*}
    \mathsf{GrandProd}
    \bigl(\ext z\bm 1+\bm w\bigr)
    =
    \mathsf{GrandProd}
    \bigl(\bm C_{\log_2 \kappa}(\ext z)\bigr).
\end{equation*}

The prover first reduces the grand product of the terminal column
\(
    \ext{\bm c}
    \coloneqq
    \bm C_{\log_2\kappa}(\ext z)
    \in\ExtensionField^m
\)
through a sequence of recursive Hadamard-product proofs. Set
\(
    \ext{\bm c}_0
    \coloneqq
    \ext{\bm c}.
\)
For each \(i=0,\ldots,\log_2m-1\), split
\(
    \ext{\bm c}_i
    =
    \ext{\bm c}_{i,L}
    \mathbin{\|}
    \ext{\bm c}_{i,R}
\)
into equal halves and define
\begin{equation*}
    \ext{\bm c}_{i+1}
    \coloneqq
    \ext{\bm c}_{i,L}
    \circ
    \ext{\bm c}_{i,R}.
\end{equation*}
Thus,
\(
    \ext{\bm c}_i
    \in
    \ExtensionField^{m/2^i},
\)
and the final layer consists of the single scalar
\(
    \ext c_{\log_2m}
    =
    \Theta
    =
    \mathsf{GrandProd}(\ext{\bm c})
    =
    \mathsf{GrandProd}
    \bigl(\ext z\bm 1+\bm w\bigr).
\)
Each transition
\[
    \ext{\bm c}_{i+1}
    =
    \ext{\bm c}_{i,L}
    \circ
    \ext{\bm c}_{i,R}
\]
is a Hadamard-product relation and is proved recursively from
\(i=\log_2m-1\) down to \(i=0\).

This produces a
tensor-structured linear claim
\begin{equation*}
    \ext c_{\log_2 \kappa}
    \EqOrNot
    \left\langle
        \ext{\bm\xi}_{\log_2 \kappa},
        \bm C_{\log_2 \kappa}(\ext z)
    \right\rangle.
\end{equation*}

The protocol then unfolds the column-product tree from
\(t=\ell_\_\kappa-1\) down to \(t=0\). Suppose that the current
claim is
\begin{equation*}
    \ext c_{t+1}
    \EqOrNot
    \left\langle
        \ext{\bm\xi}_{t+1},
        \operatorname{vec}
        \bigl(\bm C_{t+1}(\ext z)\bigr)
    \right\rangle.
\end{equation*}
By Eq.~\eqref{eq:packed-permutation-product-tree}, this is the
weighted Hadamard-product relation
\begin{equation}
\begin{aligned}
    \ext c_{t+1}
    \EqOrNot
    \Bigl\langle
        \ext{\bm\xi}_{t+1},\,
        &\operatorname{vec}
        \bigl(\bm C_{t,L}(\ext z)\bigr)
        \circ\,
        \operatorname{vec}
        \bigl(\bm C_{t,R}(\ext z)\bigr)
    \Bigr\rangle.
\end{aligned}
\label{eq:packed-permutation-hadamard-claim}
\end{equation}

Each entry of \(\bm C_t(Z)\) is represented by its base-field
coefficient vector. Its evaluation at \(\ext z\) is absorbed into
the tensor-structured weight in
Eq.~\eqref{eq:packed-permutation-hadamard-claim}. The prover and
verifier then apply the claim-reduction procedure of
Algorithm~\ref{alg:hadamard}. The two resulting MLE claims are
randomly combined into a single tensor-structured claim
\begin{equation*}
    \ext c_t
    \EqOrNot
    \left\langle
        \ext{\bm\xi}_t,
        \operatorname{vec}
        \bigl(\bm C_t(\ext z)\bigr)
    \right\rangle.
\end{equation*}
This becomes the input claim for the next level.

After \(\log_2 \kappa\) reductions, the remaining claim is
\[
    \ext c_0
    \EqOrNot
    \left\langle
        \ext{\bm\xi}_0,
        \operatorname{vec}
        \bigl(\bm C_0(\ext z)\bigr)
    \right\rangle.
\]
Since
\[
    \bm C_0(\ext z)
    =
    \ext z\bm 1_{m\times\kappa}+\bm D,
\]
this claim is equivalent to
\begin{equation*}
    \left\langle
        \ext{\bm\xi}_0,\bm d
    \right\rangle
    \EqOrNot
    \ext c_0
    -
    \ext z
    \left\langle
        \ext{\bm\xi}_0,\bm 1_N
    \right\rangle.
\end{equation*}
Thus, the proof of one grand product terminates in a single MLE
opening of the original committed vector \(\bm d\).

Applying this reduction to both \(\bm a\) and \(\bm b\), followed
by the comparison in
Eq.~\eqref{eq:packed-permutation-fingerprint}, proves the
permutation relation.

The complete protocol is given in Algorithm~\ref{alg:permutation}
in Appendix~\ref{appx:pseudocode}.

\begin{proposition}[Packed permutation proof]
\label{proposition:permutation-complexity}
The packed permutation proof is complete. Let
\(\epsilon_{\mathsf{GP}}\) denote the soundness error of proving
one grand product, excluding the terminal PCS openings. Then
\[
\begin{aligned}
    \epsilon_{\mathsf{GP}}
    \leq\;&
    \frac{
        \log_2m(\log_2m-1)
    }{
        |\ExtensionField|
    }
    +
    \frac{
        2\log_2m\log_2\kappa
    }{
        |\ExtensionField|
    }
    +
    \sum_{t=0}^{\log_2\kappa-1}
    \frac{
        2^{\log_2\kappa-t+1}-2
    }{
        |\ExtensionField|
        -
        2^{\log_2\kappa-t}
    }.
\end{aligned}
\]
Consequently, when both \(\bm a\) and \(\bm b\) are committed,
the permutation proof has soundness error at most
\[
    \min\left\{
        1,
        \frac{N-1}{|\ExtensionField|}
    \right\}
    +
    2\epsilon_{\mathsf{GP}}
    +
    \epsilon_{\mathsf{PCS}},
\]
where \(\epsilon_{\mathsf{PCS}}\) denotes the soundness error of
the terminal MLE openings.

Under schoolbook arithmetic, the prover cost for one
grand-product proof is
\[
    O(N\kappa)\,\bb
    +
    O(N)\,\be
    +
    \left(
        \frac{10N}{\kappa}
        +
        O\left(
            \kappa
            +
            (\log_2m)^2
            +
            \log_2m\log_2\kappa
        \right)
    \right)\ee.
\]
The peak prover memory is \(O(N+\kappa^2)\) base-field
coordinates, which is \(O(N)\) whenever \(m\geq\kappa\).
Proving both committed grand products doubles the corresponding
prover costs.
\end{proposition}

The proof is deferred to Appendix~\ref{appx:proof-permutation}.
